\subsection{Relación entre las vistas de arquitectura}
\label{sec:iso42010}

RT-02.03 exige cinco vistas de la arquitectura: lógica, procesos, despliegue, datos y seguridad. La Tabla~\ref{tab:vistas-iso42010} indica qué preocupación responde cada una y dónde debe demostrarse en el Subdocumento 4 o en el capítulo de datos. La integración atraviesa esas vistas; no se presenta como una sexta vista exigida por la Base.

\begin{tablalafrox}{Trazabilidad de las cinco vistas de arquitectura}{tab:vistas-iso42010}{F{0.15} F{0.34} F{0.27} Y}{\cab{Vista} & \cab{Pregunta que responde} & \cab{Lugar de la evidencia} & \cab{Responsable}}
Lógica & Quién realiza cada función y por qué contrato se comunica. & 4.1: capas, módulos, interfaces. & Arquitectura \\
Procesos & Cómo atraviesan los módulos el pedido, despacho, entrega y rendición. & 4.1: eventos y flujos; 3.4: operación. & Operaciones \\
Despliegue & Dónde se ejecuta y cómo continúa durante una falla. & 4.2 y 4.3. & Infraestructura \\
Datos & Quién posee, conserva y reconcilia cada dato. & 4.1: dominios; capítulo 5: modelo y gestión. & Datos \\
Seguridad & Quién puede actuar, con qué identidad y qué controles. & 4.1: controles lógicos; 4.2: implantación. & Seguridad \\
\end{tablalafrox}

{\footnotesize Fuente: elaboración propia de LafroX a partir de Bases Técnicas Transversales, RT-02.03.\par}

La evidencia de una vista se comprueba en el apartado indicado, no por el nombre de la tabla.

La vista lógica queda desarrollada aquí. La comprobación de las otras cuatro requiere revisar el capítulo integrado y sus figuras; esta tabla es una ruta de lectura, no una declaración automática de conformidad con ISO/IEC/IEEE 42010.

